\chapter{Estatutos}

\section{Arquitetura de equipa}

\subsection{Estatutos Pessoais}

\subsubsection{Estatutos Pessoais Executivos}
\begin{enumerate}[label=\thesubsubsection.\arabic*]
  \item Member - membro integrante da equipa.
  \item Subsystem Responsible - membro integrante da equipa, indigitado pela Direção ao Estatuto em causa.
  \item Coordinator - membro integrante da equipa, indigitado pela Board ao Estatuto em causa.
  \item Financial Manager - membro integrante da equipa, candidato eleito ao Estatuto que ocupa, ou cooptado pela Board.
  \item Department Leader - membro integrante da equipa, candidato eleito ao Estatuto que ocupa, ou cooptado pela Board.
  \item Chief Operations - membro integrante da equipa, candidato eleito ao Estatuto que ocupa, ou cooptado pela Board.
  \item Chief Engineer - membro integrante da equipa, candidato eleito ao Estatuto que ocupa, ou cooptado pela Board.
  \item Team Leader - membro integrante da equipa, candidato eleito ao Estatuto que ocupa.
\end{enumerate}

\subsubsection{Estatutos Pessoais Não Executivos}
\begin{enumerate}[label=\thesubsubsection.\arabic*]
  \item Faculty Advisor - Docente da Universidade do Porto, indigitado pela Board.
  \item Alumni - ex-membro da equipa.
\end{enumerate}

\subsection{Estatutos de Grupo}
\begin{enumerate}[label=\thesubsection.\arabic*]
  \item Comunidade - conjunto de todos as pessoas com Estatutos Pessoais supramencionados, e porventura outras que considerado pertinente.
  \item Equipa - conjunto de membros da equipa que cumpram com os Estatutos Pessoais Executivos.
  \item Departamento - conjunto de membros da equipa que se enquadram na mesma divisão de área de trabalhos arbitrada pela Direção.
  \item Direção - conjunto de membros da equipa, eleitos aos Estatutos a que se propuseram.
  \item Board - conjunto de membros da equipa, eleitos aos Estatutos a que se propuseram, com responsabilidade administrativa máxima.
\end{enumerate}

\section{Estatutos Pessoais Executivos}
\begin{enumerate}
  \item Têm um conjunto de Direitos - transversais a qualquer Estatuto Executivo:
  \begin{enumerate}
    \item Discutir sobre toda e qualquer atividade da equipa.
    \item Eleger e candidatar-se a cargos da Direção da equipa.
    \item Ter um espaço para trabalhar, tendo acesso por Cartão da Universidade do Porto ao espaço de trabalho da equipa.
    \item Candidatar-se a participar em competições.
    \item Representar a equipa em Eventos das competições, se decidido favoravelmente.
    \item Participar em Inspeções Técnicas em competições, se decidido favoravelmente.
    \item Estar presente em Dynamic Areas em competições, com o protótipo, se decidido favoravelmente.
    \item Conduzir um protótipo da equipa, se decidido favoravelmente.
    \item Efetuar cursos certificados, eventualmente obrigatórios por regras de competição.
    \item Demissão com efeito imediato, se assim o decidir.
    \item Utilização, para desenvolvimento do projeto, de Software e Hardware adquirido pela equipa.
    \item Utilizar o desenvolvimento feito no âmbito da equipa para o percurso académico, tendo aprovação da Direção.
    \item Ter acesso à documentação e outro desenvolvimento elaborado pela equipa.
    \item Ser reconhecido pelo contributo proporcionado à equipa, de um ponto de vista interno e externo.
  \end{enumerate}

  \item Têm um conjunto de Deveres - dependentes do Estatuto:
  \begin{enumerate}
    \item Aplicados ao Estatuto de Member:
    \begin{enumerate}
      \item Conhecer, respeitar e fazer respeitar o presente Regulamento Geral.
      \item Procurar compreender e seguir decisões tomadas pela Direção, sempre que esta não decida contra os seus Direitos.
      \item Manter a Direção atualizada relativamente aos seus contactos pessoais e assegurar a receção de qualquer informação enviada para os mesmos.
      \item Fornecer à Direção todos os dados pessoais estritamente necessários para o funcionamento da equipa e participação em competições.
      \item Participação em todas as reuniões de presença obrigatória: RDEP, RACP, RGER.
      \item Participação em Reuniões, se requisitado pela Direção: REXT, RISU, REMP.
      \item Participação em Eventos de Equipa e quaisquer atividades promovidas pela equipa.
      \item Fornecer à Direção o horário semanal para o semestre, até um dia util do início do mesmo.´
      \item Dedicar um conjunto mínimo de 21 horas de trabalho semanal ao projeto.
      \item Efetuar as tarefas atribuidas, procurando cumprir com a linha temporal associada.
      \item Auxiliar a equipa nas suas necessidades, de modo a cumprir com os objetivos e operações diárias.
      \item Auxiliar na criação de conteúdo para divulgação ou outros, se requisitado.
      \item Cumprir com os Valores da equipa, colocando-os acima de interesses pessoais.
      \item Basear as relações com outros membros da equipa em honestidade, confiança e comunicação.
      \item Apoiar a equipa nas suas necessidades se requisitado pela mesma, na eventualidade de uma Demissão.
      \item Desempenhar com empenho e respeito os trabalhos que forem atribuidos.
      \item Documentar o desenvolvimento efetuado, procurando transmitir conhecimento e assegurar a continuidade da equipa.
      \item Defender o nome e interesses da FS FEUP.
      \item Assegurar a manutenção da confidencialidade de temas da equipa cuja exposição dentro ou fora da mesma seja considerado impertinente.
    \end{enumerate}

    \item Aplicados ao Estatuto de Subsystem Responsible:
    \begin{enumerate}
      \item Os Deveres supramencionados associados ao Estatuto de Member.
      \item Promover e defender os interesses e requisitos do respetivo Subsystem.
      \item Ser o ponto primário de contacto entre o Department Leader e o Subsystem.
      \item Auxiliar na coordenação dos demais membros do Subsystem.
      \item Ser responsável por decisões elaboradas no decorrer do desenvolvimento do seu Subsystem.
    \end{enumerate}

    \item Aplicados ao Estatuto de Coordinator:
    \begin{enumerate}
      \item Os Deveres supramencionados associados ao Estatuto de Member.
      \item Cumprir com empenho as funções específicas segundo as quais foi indigitado pela Board.
    \end{enumerate}

    \item Aplicados ao Estatuto de Financial Manager:
    \begin{enumerate}
      \item Os Deveres supramencionados associados ao Estatuto de Member.
      \item Revisão do Planeamento e Orçamento Logístico.
      \item Elaboração de estratégias de perceção orçamental da equipa.
      \item Gestão orçamental e financeira do projeto.
      \item Elaboração e atualização do Balancete da equipa, bem como a sua coerência.
      \item Elaboração dos Documentos de Relatórios e Contas/Orçamento.
      \item Preparação e detalhe das transações da equipa.
      \item Pós-processamento dos documentos associados.
      \item Cobrir quaisquer responsabilidades que possam advir de uma lacuna em termos de desenvolvimento financeiro do projeto.
    \end{enumerate}

    \item Aplicados ao Estatuto de Department Leader:
    \begin{enumerate}
      \item Os Deveres supramencionados associados ao Estatuto de Member.
      \item Coordenar de um ponto de vista técnico e pessoal os membros que formam o respetivo Departamento.
      \item Responsabilidade de validação do trabalho feito pelos membros do Departamento.
      \item Ferramenta de validação da documentação elaborada intradepartamento, assegurando coerência e rigor na mesma.
      \item Gestão do trabalho, de modo ao Departamento cumprir com a linha temporal e objetivos definidos.
      \item Promoção do bem estar e motivação dos membros integrantes do respetivo Departamento.
      \item Presença ou delegação da mesma, no que concerne a REMP ou contactos relacionados com necessidades do respetivo Departamento.
      \item Ser o pronto primário de contacto entre a Direção e Board e o Departamento.
      \item Auxílio na gestão orçamental do respetivo Departamento.
      \item Cobrir quaisquer responsabilidades que possam advir de uma lacuna intradepartamento em termos de desenvolvimento.
    \end{enumerate}

    \item Aplicados ao Estatuto de Chief Operations:
    \begin{enumerate}
      \item Os Deveres supramencionados associados ao Estatuto de Member.
      \item Coordenar de um ponto de vista técnico e pessoal os membros que formam o respetivo Departamento.
      \item Responsabilidade de validação do trabalho feito pelos membros do Departamento.
      \item Ferramenta de validação da documentação elaborada intradepartamento, assegurando coerência e rigor na mesma.
      \item Gestão do trabalho, de modo ao Departamento cumprir com a linha temporal e objetivos definidos segundo a Direção.
      \item Presença ou delegação da mesma, no que concerne a REMP ou contactos relacionados com necessidades do respetivo Departamento.
      \item Ser o pronto primário de contacto entre a Direção e Board e o Departamento.
      \item Auxílio na gestão orçamental do respetivo Departamento.
      \item Responsabilidade pela tutela dos Static Events de Cost and Manufacturing e Business Plan Presentation.
      \item Assegurar a aquisição de patrocínios necessários à execução do projeto.
      \item Coordenação das logísticas associadas à execução do projeto: competições, Rollout, Eventos de equipa, Quizzes, e qualquer outra.
      \item Promoção do bem estar e motivação dos membros integrantes do respetivo Departamento e equipa.
      \item Responsabilidade pela organização de eventos que concernem a equipa.
      \item Promoção e responsabilidade pela imagem externa e interna da equipa, em termos de conteúdo visual produzido.
      \item Cobrir quaisquer responsabilidades que possam advir de uma lacuna intradepartamento, ou do projeto, em termos de desenvolvimento.
    \end{enumerate}

    \item Aplicados ao Estatuto de Chief Engineer:
    \begin{enumerate}
      \item Os Deveres supramencionados associados ao Estatuto de Member.
      \item Coordenação técnica e pessoal dos diversos Departamentos da equipa cuja área do saber envolve conhecimentos de Engenharia Mecânica, Engenharia Eletrotécnica e de Computadores, Informática e demais semelhantes.
      \item Responsabilidade de defender o desenvolvimento do melhor produto possível, sempre consciente do orçamento e planeamento da equipa.
      \item Ferramenta de revisão global de documentação técnica produzida pela equipa.
      \item Responsabilidade de coordenar as Technical Inspections e seu planeamento, assegurando que a equipa as cumpre.
      \item Tomar decisões técnicas nas interfaces entre Subsystems, procurando a obtenção de um produto final rigoroso e coerente.
      \item Responsabilidade de coordenar as fases de testagem dos Subsystems e do protótipo completo.
      \item Responsável máximo, em testagem e competição, por validar que o protótipo se encontra seguro e passível de ser testado, de um ponto de vista mecânico.
      \item Responsabilidade pela tutela do Engineering Design Event.
      \item Cobrir quaisquer responsabilidades que possam advir de uma lacuna em termos de desenvolvimento do projeto.
    \end{enumerate}

    \item Aplicados ao Estatuto de Team Leader:
    \begin{enumerate}
      \item Os Deveres supramencionados associados ao Estatuto de Member.
      \item Ser o ponto de contacto principal entre os parceiros institucionais e a equipa, em colaboração com o Faculty Advisor.\%: Reitoria da Universidade do Porto, Direção da FEUP, Direção dos Departamentos da FEUP, Direção da AEFEUP, Direção da FAP, e demais Instituições e Núcleos.
      \item Ser o ponto de contacto principal entre a presente equipa e Comunidade FS FEUP.
      \item Definição do orçamento global para o ciclo de equipa em causa.
      \item Estabelecimento do planeamento a longo prazo da equipa, que se reflete no planeamento macro do ciclo.
      \item Promoção do bem estar e motivação dos membros integrantes da equipa.
      \item Promover ações que levem a um bom funcionamento da equipa.
      \item Coordenação e curadoria do processo de Recrutamento da equipa.
      \item Assegurar que a equipa possui dos recursos necessários à execução do projeto.
      \item Assegurar que o máximo potencial é passível de ser extraído da equipa.
      \item Defender os interesses da equipa e do coletivo dos seus membros, acima de tudo, de um ponto de vista interno e externo à equipa.
      \item Cobrir quaisquer responsabilidades que possam advir de uma lacuna em termos de desenvolvimento do projeto.
      \item Ter a capacidade de substituir em funções, se necessário, qualquer membro da equipa.
      \item Promover decisões que procurem ir de encontro à Visão da equipa.
    \end{enumerate}
  \end{enumerate}
\end{enumerate}

\section{Estatutos Pessoais Não Executivos}
\begin{enumerate}
  \item Englobam os Estatutos de:
  \begin{enumerate}
    \item Faculty Advisor.
    \item Alumni.
  \end{enumerate}
  \item O Estatuto de Faculty Advisor tem um conjunto de Direitos:
  \begin{enumerate}
    \item Opinar e discutir sobre toda e qualquer atividade da equipa.
    \item Participar em competições com a equipa.
    \item Efetuar cursos certificados, eventualmente obrigatórios por regras de competição.
    \item Demissão com efeito imediato, se assim o decidir.
    \item Ter acesso a todo a documentação e outro desenvolvimento elaborado pela equipa.
    \item Ser reconhecido pelo contributo proporcionado à equipa, de um ponto de vista interno e externo.
  \end{enumerate}
  \item O Estatuto de Faculty Advisor tem um conjunto de Deveres:
  \begin{enumerate}
    \item Promover o bem estar da totalidade dos membros que integram a equipa.
    \item Ser o ponto de contacto entre os parceiros institucionais e a equipa, em colaboração com o Team Leader.
    \item Zelar pelo cumprimento do presente Regulamento Geral e demais legislação aplicada aos membros da equipa.
    \item Assegurar uma adequada transição de ciclo de equipa e continuidade, em coordenação com a Direção em final de Mandato.
    \item Acompanhamento da equipa em competições, assegurando a segurança e bem estar dos membros da equipa, assim como o cumprimento das regras de conduta prescrita pelo país da competição e regras da competição.
    \item Contribuir como ferramenta de validação do trabalho desenvolvido pela equipa, de acordo com a sua área de saber enquanto docente na Universidade do Porto.
  \end{enumerate}
  \item O estatuto de Alumni tem um conjunto de Direitos:
  \begin{enumerate}
    \item Opinar e discutir sobre toda e qualquer atividade da equipa.
    \item Ser reconhecido pelo contributo proporcionado à equipa, de um ponto de vista interno e externo.
    \item Acompanhar o desenvolvimento feito pela equipa.
    \item Participar em Eventos de Equipa.
  \end{enumerate}
  \item O estatuto de Alumni tem um conjunto de Deveres:
  \begin{enumerate}
    \item Intervir em processos de validação do desenvolvimento feito pela equipa, se disponível e assim requisitado.
    \item Zelar pelo cumprimento do presente Regulamento Geral e demais legislação aplicada aos membros da equipa.
    \item Promover o bem estar da totalidade dos membros que integram a equipa.
    \item Assegurar uma adequada transição de ciclo de equipa e continuidade, em coordenação com a Direção em final de Mandato.
    \item Criar um ambiente construtivo de melhoria contínua da equipa e protótipos, contribuindo para a criação de um ambiente propício ao desenvolvimento da equipa.
  \end{enumerate}
\end{enumerate}

\section{Estatutos de Grupo}
\begin{enumerate}
  \item Responsabilidades da Comunidade:
  \begin{enumerate}
    \item Ser composta por pessoas de todos os Estatutos, desde que tenham como objetivo promover o bem estar da totalidade dos membros que integram a equipa.
    \item Zelar pelo cumprimento do presente Regulamento Geral e demais legislação aplicada aos membros da equipa.
    \item Assegurar uma adequada transição de ciclo de equipa e continuidade, em coordenação com a Direção em final de Mandato.
    \item Criar um ambiente construtivo de melhoria contínua da equipa e protótipos, ultrapassando interesses pessoais que possam entrar em conflito com o primeiro.
    \item Defender os feitos do passado da equipa como meio para proporcionar um melhor futuro à mesma.
  \end{enumerate}
  \item Responsabilidades da Equipa:
  \begin{enumerate}
    \item Entreajudar-se no desenvolvimento pessoal individual e objetivos coletivos.
    \item Zelar pelo cumprimento do presente Regulamento Geral e demais legislação aplicada aos membros da equipa.
    \item Defender e cumprir diariamente com os Valores da equipa.
    \item Procurar cumprir com os objetivos definidos para o projeto no ciclo.
  \end{enumerate}
  \item Responsabilidades de Departamento:
  \begin{enumerate}
    \item Cumprir a demais legislação aplicada ao Estatuto de Equipa.
    \item Promover os interesses do Departamento em que se incluem, não descurando uma perspetiva global do enquadramento do Departamento na equipa e seus objetivos.
    \item Promover um ambiente de confiança na liderança do Departamento, não obstante escrutínio construtivo das decisões e metodologia seguidas.
  \end{enumerate}
  \item Responsabilidades da Direção:
  \begin{enumerate}
    \item Cumprir a demais legislação aplicada ao Estatuto de Equipa.
    \item Ser composta por um número mínimo de seis membros da equipa: Team Leader, Chief Engineer e quatro Department Leaders.
    \item Representar a equipa perante os parceiros institucionais, empresariais e outros.
    \item Coordenar os trabalhos a nível interno e externo.
    \item Coordenar os membros da equipa.
    \item Assumir responsabilidade pelas ações dos membros da equipa, no âmbito da atividade da mesma.
    \item Gerir o património da equipa, a nível de manutenção ou cedência do mesmo.
    \item Garantir o correto funcionamento da equipa, assumindo as responsabilidades necessárias para tal.
    \item Zelar pelo cumprimento do presente Regulamento Geral e demais legislação aplicada aos membros da equipa.
    \item Defender a equipa e seus interesses acima de qualquer outra motivação.
  \end{enumerate}
  \item Responsabilidades da Board:
  \begin{enumerate}
    \item Cumprir a demais legislação aplicada ao Estatuto de Equipa e Direção.
    \item Ser composta por um número mínimo de dois membros da equipa: Team Leader e Chief Engineer.
    \item Gerir os fluxos financeiros da equipa.
    \item Coordenar os membros da Direção.
    \item Tomar decisões de âmbito abrangente em termos temporais e de gestão da equipa.
    \item Assumir responsabilidade pelas decisões tomadas pela equipa e suas consequências.
    \item Curadoria de diversos processos associados à manutenção e continuidade da equipa.
    \item Garantir que o coletivo dos Deveres da Board é cumprido, independentemente da constituição de Estatutos da Board.
    \item Indigitar o Faculty Advisor.
  \end{enumerate}
\end{enumerate}
